\chapter*{Введение}
\addcontentsline{toc}{chapter}{Введение}

Прогнозирование состояния атмосферы и связанных с ним гидрометеорологических процессов относится к числу наиболее значимых прикладных задач современной науки о данных. От точности прогноза зависят системы раннего предупреждения опасных явлений, управление транспортной инфраструктурой, оценка рисков для энергетики и сельского хозяйства, а также принятие решений в условиях неопределенности. Наиболее сложными для прогноза являются осадки: они обладают выраженной пространственно-временной неоднородностью, сильной асимметрией распределения, большим числом нулевых наблюдений и редкими, но важными экстремальными выбросами.

В последние годы в задачах прогноза погоды активно развиваются нейросетевые подходы, обучаемые на крупных архивах реанализа и наблюдений. Архитектуры GraphCast, Pangu-Weather, FuXi и близкие к ним модели показали, что глубокие модели могут быть конкурентоспособны по сравнению с традиционными системами среднесрочного прогноза \cite{graphcast2023,pangu2023,fuxi2023}. При этом практическое использование таких моделей требует аккуратной постановки эксперимента: необходимо определить тип данных, способ задания редкого события, набор метрик качества и базовые модели для сравнения.

В работе рассматриваются две взаимодополняющие постановки. Первая связана со станционными временными рядами и прогнозом суммарных осадков на горизонтах от 24 до 240 часов. Для нее используются табличное представление данных, признаки на основе лагов и скользящих статистик, а также базовые модели различной сложности. Вторая постановка использует сеточные региональные поля и ориентирована не только на прогноз количества осадков, но и на вероятностную детекцию экстремальных осадочных событий.

\textbf{Целью работы} является разработка и исследование подходов к прогнозированию осадков и детекции аномальных осадочных событий на основе станционных и сеточных климатических данных с использованием методов машинного обучения и глубоких нейросетевых моделей.

Для достижения поставленной цели в работе решаются следующие \textbf{задачи}:
\begin{enumerate}[label=\arabic*)]
    \item сформулировать задачу детекции аномальных осадков через превышение высококвантильного порога;
    \item подготовить признаки для станционного прогноза суммарных осадков;
    \item реализовать и сопоставить климатологию, линейную модель и градиентные бустинги как базовые методы;
    \item описать региональную нейросетевую модель для многогоризонтного прогноза осадков на сеточных данных;
    \item оценить вероятностные характеристики прогноза экстремальных осадков;
    \item проанализировать ошибки модели по горизонту, пространству, времени и уровню неопределенности.
\end{enumerate}

\textbf{Объектом исследования} являются климатические и метеорологические данные, описывающие осадки на станционном и сеточном уровнях.

\textbf{Предметом исследования} являются методы машинного обучения и глубокого обучения для прогнозирования осадков и детекции аномальных осадочных событий.

\textbf{Научная новизна} состоит в адаптации единой экспериментальной схемы для анализа осадков как непрерывной величины и как редкого события превышения локального квантильного порога. В работе отдельно рассматриваются сильные стороны станционных табличных моделей и возможности региональной нейросетевой модели, работающей с пространственными полями.

\textbf{Практическая значимость} связана с тем, что рассмотренные подходы могут использоваться при построении систем краткосрочного и среднесрочного прогноза осадков, а также при оценке риска экстремальных гидроклиматических событий.

Структура работы следующая. В первой главе формулируется постановка задачи и вводятся критерии качества. Во второй главе рассматривается станционная постановка: данные, признаки и базовые модели. В третьей главе описывается сеточная постановка и архитектура региональной нейросетевой модели. В четвертой главе приводятся результаты экспериментов, графики, калибровка и анализ ошибок. Дополнительные конфигурационные таблицы встроены в основной текст.
